% RASCUNHO PRELIMINAR redigido com apoio de IA (Claude Code, Anthropic; ver ficha de uso de IA).
% A interpretação é uma proposta de leitura dos dados e deve ser revisada, validada e reescrita pelo autor.
\section{Organização do Trabalho}

Este trabalho está organizado em cinco capítulos. O Capítulo~\ref{chp:fund_teo} apresenta a fundamentação teórica: a virtualização de redes, as redes definidas por software, a virtualização de funções de rede, o fatiamento de redes, a demanda elástica e os algoritmos de alocação de fatias considerados (PETIC, ACO, SA e Q-Learning). O Capítulo~\ref{chp:metodologia} descreve a metodologia: a caracterização da pesquisa, o ambiente de simulação, o modelo do problema, a implementação e a calibração dos algoritmos, as topologias, os cenários de demanda, as métricas e o procedimento experimental. O Capítulo~\ref{chp:resultados} apresenta e discute os resultados dos experimentos. Por fim, o Capítulo~\ref{chp:conclusao} apresenta as conclusões, as limitações e os trabalhos futuros.
